% Cote 119: offprints, reports, courses, CV and letters (typed and printed part): archive cote 119, site page 286 (PDF page 287).
% typed letter. Draft compared: none.
% AI-transcribed by Claude Opus 5.5 (claude-opus-5-5, adaptive thinking, effort high), Message Batches API, batch r119_full (msgbatch_01DwZNm5Y7YsKBWyAN4sFs8w), 2026-11-13; not yet checked by eye.
\scanpage{119/286}
problems that had not been thought of before, that were technically fairly
easy to solve, and for which the answers were important. I was looked at as
if I was asking for the philosopher's stone! Or even that it might be
easier to find some standard problem and solve that. Afterwards, I thought
of an old good story, which you may know, but which is good for a laugh if
not, and is frighteningly pertinent. A drunk is searching on the ground
under a lamppost at night. A man asks him what is the matter. `I am
looking for my watch.' `Did you lose it here?' `No! I lost it over
there.' `Why are you looking here?' `The light is better here!'

The point is that a lot of people cannot understand why I am looking at
such `peculiar' (by their lights) areas of mathematics. But what is the
point in looking hard at the areas which everyone accepts as properly
formulated? French mathematicians understand this point much better than UK
Mathematicians. It is also something to do with the training. So one can
hope in a small way to influence some people at Bangor in a more open
approach. Young people are starved of a proper discussion of methodology,
of the practice of research mathematics, of a global point of view. This
difficulty is widespread. I speak to people from Egypt, Jordan, UK,..., and
the story is the same. Students ask: why are we doing these courses? and
the only answer for which the lecturers are able to give is: To pass the
examinations! It is all rather sad.

I have contact with Muster Coulibaly. Do you know him? I write
references for him, but do not know much in detail about his work.

\hand{Ronnie} \hfill \hand{Yours affectionately}

\hfill \hand{Ronnie}

\mspage{5}

\margin{286}
